% Responsable: Arteaga Alava Danela Dayana · Ronda 2
\section{Introducción y dominio}

Este trabajo aplica tres técnicas de priorización de requisitos a un sistema de gestión para un gimnasio. Las técnicas son MoSCoW, el modelo de Kano y la votación ponderada de 100 puntos.

El sistema propuesto cubre tres procesos del gimnasio. El primero es la gestión de membresías, que incluye el registro de socios, la contratación de planes, los pagos y los vencimientos. El segundo es el control de asistencia en recepción. El tercero es la asignación y consulta de rutinas de entrenamiento. Las partes interesadas consideradas son el administrador o propietario, el personal de recepción, los entrenadores y los socios.

El propósito es comprobar que la prioridad de un requisito depende de quién la decide y de la técnica que se usa. Para ello se redactan 20 requisitos según ISO/IEC/IEEE 29148 \cite{ieee29148}, se clasifican con MoSCoW, se aplica el cuestionario de Kano a 8 características y se realiza una votación de 100 puntos. Después se comparan los tres órdenes con la correlación de rangos de Spearman y se propone la lista de requisitos para una primera versión del sistema.
